
\section{Scale-invariant torque, angular-momentum compensation and smooth packet stacks}
\label{sec:xxxviii-scale}
The exact nonaxisymmetric zero-force Gaussian vortex of Edition~XXXVII has
\[
J(r,z)=\overline{u'_r\Gamma'}=\varepsilon^2r^3e^{-2(r^2+z^2)},\qquad
\mathcal T=-r^{-1}\partial_r(rJ)
=4\varepsilon^2r^2(r^2-1)e^{-2(r^2+z^2)}.
\]
Here $\Gamma=ru_\theta$ and $\varepsilon=\nu/1000$.
Although its initial torque is nonzero, it redistributes angular momentum
without generating a nonzero total.

\begin{proposition}[Exact interior--exterior angular-momentum compensation]
\label{prop:xxxviii-compensation}
For each $R>0$ and $z$,
\begin{equation}\label{eq:xxxviii-balance}
\boxed{\quad\int_0^R\mathcal T(0,r,z)\,r\,dr
=-\varepsilon^2R^4e^{-2(R^2+z^2)}.\quad}
\end{equation}
In particular, $\int_0^\infty\mathcal T(0,r,z)r\,dr=0$.
For any sufficiently decaying smooth unforced solution,
the global $z$-angular momentum
$\int_{\mathbb R^3}(xu_y-yu_x)\,dx$ is conserved.
\end{proposition}
\begin{proof}
Integrate $r\mathcal T=-\partial_r(rJ)$;
the boundary term at $r=0$ vanishes as $r^4$ and the
infinite boundary term is zero by Gaussian decay.
Testing the full physical Navier--Stokes momentum equation
against the rotational Killing field $(-y,x,0)$ proves global
angular-momentum conservation by the symmetry of the nonlinear
and viscous stresses. The pressure contributes no torque.
\end{proof}

\begin{theorem}[Arbitrarily large instantaneous torque without singularity]
\label{thm:xxxviii-scaling}
Let $u$ be the globally smooth zero-force
solution of Edition~XXXVII and define the exact
Navier--Stokes rescaling
$u_\lambda(x,t)=\lambda u(\lambda x,\lambda^2t)$,
$\lambda>0$. Then every $u_\lambda$ is itself
globally smooth and unforced, and
\[
\boxed{\quad
\Gamma_\lambda(x,t)=\Gamma(\lambda x,\lambda^2t),
\qquad
\mathcal T_\lambda(x,t)=\lambda^2\mathcal T(\lambda x,\lambda^2t).
\quad}
\]
At its rescaled radial sample,
\begin{equation}\label{eq:xxxviii-instantaneous}
\boxed{\quad
\mathcal T_\lambda(0,1/(2\lambda),0)
=-\frac34\varepsilon^2\lambda^2e^{-1/2}.\quad}
\end{equation}
Nevertheless
\[
\boxed{\quad
\|u_\lambda(0)\|_2^2=\lambda^{-1}\|u(0)\|_2^2,
\qquad
\|u_\lambda(0)\|_{\mathcal X^{-1}}
=\|u(0)\|_{\mathcal X^{-1}}.\quad}
\]
For each fixed $\lambda$, the time-integrated
$\|\mathcal T_\lambda\|_\infty$ is bounded up to
$t=1$ and angular variation on the shrinking source
circle $r_\tau=O(\sqrt\tau)$ is $O_\lambda(\sqrt\tau)$.
These bounds need not be uniform in $\lambda$.
\end{theorem}
\begin{proof}
Directly substitute the parabolic scaling into the
unforced PDE, cylindrical angular-momentum definition
and radial/axial covariance flux. The unitary
Fourier scaling
$\widehat{u_\lambda}(\xi)=\lambda^{-2}
\widehat u(\xi/\lambda)$ proves the critical
$\mathcal X^{-1}$ invariance. Change variables
in the spatial $L^2$ norm for the kinetic-energy relation.
The sample torque follows from the exact Gaussian formula.
Fixed-scale global smoothness supplies bounded derivatives
on the compact interval $[0,1]$ and hence the finite
torque integral. The identity
$\Gamma=xu_y-yu_x$ gives
$|\partial_\theta\Gamma|
\le r\|u\|_\infty+r^2\|\nabla u\|_\infty$,
and proves the angular-variation assertion.
Letting $\lambda$ vary with $\tau$ does not define
a single initial-value evolution.
\end{proof}

\begin{theorem}[Smoothness obstruction for a direct infinite packet cascade]
\label{thm:xxxviii-packets}
Choose a nonzero solenoidal compact field
$\phi=\nabla\times[\chi e^{-\rho^2}(x+(x^2+y^2)y)e_z]$
supported in the unit ball, with $\chi=1$
near $(r,z)=(1/2,0)$. Take
$\lambda_n=16^n$, $c_n=(0,0,4^{-n})$ and
$U_n(x)=a_n\lambda_n\phi(\lambda_n(x-c_n))$.
The packet supports are disjoint and accumulate at the origin.
Then $U_0=\sum_{n\ge1}U_n$ extends as a
$C^\infty$ velocity across the origin if and only if
\begin{equation}\label{eq:xxxviii-smoothness}
\boxed{\quad
|a_n|\lambda_n^{k+1}\to0
\qquad\text{for every fixed integer }k\ge0.\quad}
\end{equation}
At the packet sample $(r,z)=(1/(2\lambda_n),4^{-n})$
the exact initial torque is
\begin{equation}\label{eq:xxxviii-stack-torque}
\boxed{\quad
\mathcal T_{U_0}(0,1/(2\lambda_n),4^{-n})
=-\frac34 a_n^2\lambda_n^2e^{-1/2}
=o(\lambda_n^{-2}).\quad}
\end{equation}
Thus fixed critical-amplitude copies or coefficients
decaying by one fixed power of $\lambda_n$ do
not provide smooth initial data for this direct stack.
\end{theorem}
\begin{proof}
The support radii $\lambda_n^{-1}=16^{-n}$
are smaller than the gaps between the centers $4^{-n}$.
On the $n$th packet the order-$k$ derivative supremum is
$|a_n|\lambda_n^{k+1}\|\nabla^k\phi\|_\infty$.
For every fixed $k$ this derivative norm of a nonzero
compact smooth field is positive, since otherwise it
would be a compactly supported polynomial.
Derivatives of any smooth extension must vanish at the
accumulation point, as they vanish on nearby support gaps.
Continuity proves necessity of the displayed jet condition;
the same decay proves sufficiency by extension through zero.
On the local plateau the quadratic Reynolds torque
scales by $a_n^2\lambda_n^2$; other packets vanish there.
Finally the $k=1$ condition yields
$a_n\lambda_n^2\to0$, so
$a_n^2\lambda_n^2=o(\lambda_n^{-2})$.
\end{proof}


\begin{proposition}[A finite-energy but nonsmooth infinite packet stack]
\label{prop:xxxviii-finite-energy}
In Theorem~\ref{thm:xxxviii-packets} put $a_n=1$ for all $n$.
The resulting solenoidal field is in $L^2(\mathbb R^3)$
with the exact norm
\begin{equation}\label{eq:xxxviii-energy}
\boxed{\quad\|U_0\|_2^2=
\|\phi\|_2^2\sum_{n\ge1}\lambda_n^{-1}
=\frac1{15}\|\phi\|_2^2<\infty.\quad}
\end{equation}
Nevertheless it is unbounded near the accumulation point,
and therefore does not define admissible smooth initial data.
Its initial packet torques can grow as $\lambda_n^2$
only because the smoothness condition is violated.
\end{proposition}
\begin{proof}
The supports are disjoint and the three-dimensional
rescaling gives
$\|U_n\|_2^2=\lambda_n^{-1}\|\phi\|_2^2$.
Sum the geometric series $\sum_{n\ge1}16^{-n}=1/15$.
On the $n$th support,
$\|U_n\|_\infty=\lambda_n\|\phi\|_\infty\to\infty$,
with nonzero maxima at points approaching the origin.
\end{proof}

\paragraph{Near-axis consequence and the remaining open dynamics.}
For uniformly $C^4$ Cartesian velocity fields near the axis,
angular Taylor expansion and parity imply
$\overline{u'_r\Gamma'}=O(r^3)$,
$\overline{u'_z\Gamma'}=O(r^2)$
and $\mathcal T=O(r^2)$.
These constants need not remain bounded near an
unknown singular time. The three theorems rule out
large \emph{instantaneous} torque, global angular-momentum
redistribution and the naive initial packet stack
as substitutes for nonlinear terminal amplification.
They do not bound dynamically transferred energy into
smaller scales at later times. The original OpenAI
Navier--Stokes source includes forcing, and no smooth-data
zero-force singularity or full Lean PDE proof
is established by this edition.
